\documentclass[tikz]{standalone}
\usepackage{amsmath}
\usetikzlibrary{calc}
\definecolor{bg}{HTML}{171D2B}
\definecolor{aksen}{HTML}{F7B538}
\definecolor{teks}{HTML}{FFFFFF}
\begin{document}
\pagecolor{bg}
\begin{tikzpicture}[>=stealth, line width=1.1pt]
  % koordinat dasar
  \coordinate (C) at (0,0);
  \def\r{4}
  % lingkaran utama
  \draw[teks] (C) circle (\r);
  % titik pusat
  \fill[aksen] (C) circle (2.5pt) node[below left=1pt, teks] {$C$};
  % titik singgung T di (r,0)
  \coordinate (T) at (\r,0);
  \fill[aksen] (T) circle (2.5pt);
  % titik A di sudut 55 derajat
  \coordinate (A) at (55:\r);
  \fill[aksen] (A) circle (2.5pt) node[above left=2pt, teks] {$A$};
  % jari-jari ke A dan ke T
  \draw[teks] (C) -- (A);
  \draw[teks] (C) -- (T);
  % tali AB dan busur AB
  \coordinate (B) at (155:\r);
  \draw[aksen, line width=1.6pt] (A) -- (B);
  \fill[aksen] (B) circle (2.5pt) node[above left=2pt, teks] {$B$};
  \draw[teks!40, dashed, line width=0.8pt] (A) arc (55:155:\r);
  % sudut pusat alpha
  \draw[aksen, ->] (25:1.1) arc (25:55:1.1);
  \node[teks] at (40:1.5) {$\alpha$};
  % garis singgung di T
  \draw[aksen] ($(T)+(0.9,0.9)$) -- (T) -- ($(T)+(0.9,-0.9)$);
  \node[teks] at (6.2, 1.4) {garis singgung};
  % label jari-jari dan tali
  \node[teks] at (27.5:2.55) {$r$};
  \node[teks] at (105:4.55) {tali $AB$};
  % sudut antara tali dan garis singgung di A
  \draw[aksen, ->] ($(A)+(0:0.85)$) arc (0:-32.5:0.85);
  \node[teks] at ($(A)+(-0.05,-1.2)$) {$\tfrac{1}{2}\alpha$};
\end{tikzpicture}
\end{document}
